% ======================================================================================
%  Tao 2019 — "Almost all orbits of the Collatz map attain almost bounded values"
%  Lean 4 formalization blueprint (arXiv:1909.03562v5).
%
%  This is the reviewer-facing proof map: LaTeX narrative + `\lean{}` bindings (to the
%  canonical declaration names fixed in SKELETON-SPEC.md) + `\uses{}` dependency edges
%  (following the "depends" column of BLUEPRINT.md's node ledger) + `\leanok`/`\notready`
%  status.  Node ids are the labels: S1--S4 (support), C1--C10 (core spine), X1--X11
%  (the §7 crux).  The prose statement of each node is transcribed from BLUEPRINT.md's
%  descriptions and SKELETON-SPEC.md's Lean statements; where the paper equation is not
%  reproduced verbatim here a `% TODO-PRECISE` marks the coarse spot.
%
%  STATUS DISCIPLINE (updated 2026-07-10, post judge pass):
%    * Statement-`\leanok` is set on the 19 nodes whose `\lean{}`-bound statements are
%      landed, compiled, and judge-ratified (S1--S4 C1--C7 C9 C10 X2 X3 X4 X6 X7 X11).
%      The 6 nodes with no canonical Lean name yet (C8 X1 X5 X8 X9 X10) keep `\notready` until
%      their statements are pinned.  PROOF-`\leanok` (green background) is set ONLY on a
%      dated judge-run `#print axioms` verification of the bound decls -- the manual form
%      of the future BlueprintAudit reconciler; never on a worker's claim alone.
%      (X3 verified 2026-07-10: black_structure = [propext, Classical.choice, Quot.sound].)
%    * Nodes with a canonical Lean name in SKELETON-SPEC.md carry `\lean{...}`.  Nodes the
%      spec leaves un-named (e.g. Lemma 2.2, Prop 5.2, several §7 case lemmas) carry no
%      `\lean{}` and are described in prose; they will be bound once named.
%    * Every node whose PROOF is open carries a `\lapsrisk{laps}{risk}{conf}` campaign
%      estimate mirroring BLUEPRINT.md §2's ledger (laps = expected Opus-treadmill laps,
%      risk word from the as-stated confidence: >=85 low, 75-80 medium, <=70 high).
%      Drop the line only when the node's proof is done; re-sync numbers from the ledger.
% ======================================================================================

% Campaign estimate annotation (one per incomplete node). Defined here so both the
% plasTeX web build (web.tex) and the pdf build (print.tex) pick it up via \input.
\newcommand{\lapsrisk}[3]{\par\smallskip\noindent\textit{\textbf{Campaign estimate:}
#1 treadmill laps; risk \textbf{#2} (#3 confidence the node completes as stated).}}

\chapter{Program overview}

This is a blueprint for a \emph{full Lean~4 formalization of Terence Tao's 2019 theorem}
(\texttt{arXiv:1909.03562v5}) that almost all Collatz orbits attain almost bounded values.
The summit is \textbf{Theorem~1.3}: for every \(f:\mathbb{N}\to\mathbb{R}\) with
\(f(N)\to\infty\), one has \(\mathrm{Colmin}(N) < f(N)\) for almost all \(N\) in the sense of
\emph{logarithmic density}, where \(\mathrm{Colmin}(N)\) is the minimum of the Collatz orbit
of \(N\).  As of this writing the theorem is formalized in \emph{no} proof assistant, by any
route; a Lean proof is a first.

\section{The spine}

The whole argument is a chain of reductions.  Read right-to-left, each arrow is
``is proved from'':
\[
  \underbrace{\text{Thm 1.3}}_{\text{main}}
  \Longleftarrow \text{Thm 1.6}
  \Longleftarrow \underbrace{\text{Thm 3.1}}_{\text{quantitative}}
  \Longleftarrow \underbrace{\text{Prop 1.11}}_{\text{stabilization}}
  \Longleftarrow \big[\,\text{Prop 1.9} + \text{Prop 1.14}\,\big]
  \Longleftarrow \underbrace{\text{Prop 1.17}}_{\text{character decay}}
  \Longleftarrow \text{Prop 7.1}
  \Longleftarrow \text{Prop 7.3}
  \Longleftarrow \text{Prop 7.8}.
\]
The passage from the analytic main theorem to a purely arithmetic statement about the
\emph{Syracuse} map (the odd-to-odd accelerated Collatz map) happens at
Thm~1.6/Thm~3.1 (node~\ref{C6}), and drives everything down to a single Fourier-analytic
input, the \textbf{character-sum decay} Prop~1.17 (node~\ref{X11}).  Prop~1.17 is proved in
§7 by a renewal-process argument played against the geometry of the ``black'' set of the
phase functions --- the deepest part of the paper.

\section{The three risk kernels}

Almost all of the completion risk concentrates in three clusters; everything outside them is
standard treadmill fare (counting, harmonic sums, \(\mathbb{Z}/3^n\mathbb{Z}\) arithmetic,
PMF calculus).
\begin{itemize}
  \item \textbf{S3 --- the local 2-D bound} (Lemma~2.2, node~\ref{S3}): Gaussian-type local
        limit upper bounds for sums of the lattice distributions.  Classical analysis;
        long but self-contained via explicit formulas and exponential tilting (D5).
  \item \textbf{X3 --- the triangle structure} (Lemma~7.4, node~\ref{X3}): the black set of
        the phase function is a disjoint union of well-separated triangles.  Finite,
        elementary, delicate case analysis.
  \item \textbf{X8/X10 --- renewal versus triangles} (Case~2 and Lemma~7.10,
        nodes~\ref{X8},~\ref{X10}): the paper's pinnacle.  The design decision D6 finitizes
        the infinite renewal process into an explicit downward recursion \(Q\), turning the
        stopping-time measure theory into strong inductions over that recursion.
\end{itemize}
The critical path is
\(\text{S3}\to\text{X6}\to\{\text{X8},\text{X10}\}\to\text{X11}\to\text{C10}\to\text{C9}
  \to\text{C6}\).

\section{Global modelling decisions}

The Lean development departs from the paper's presentation in eight deliberate ways
(full rationale in \texttt{BLUEPRINT.md} §0); the ones that shape the statements below are:
\begin{itemize}
  \item[\textbf{D1}] \emph{Probability is discrete.}  Every random variable is a \texttt{PMF}
        on a countable type; expectations are \texttt{tsum}s and total variation is
        \(d_{\mathrm{TV}}(p,q)=\sum_a |p(a)-q(a)|\).  No measure theory (node~\ref{S1}).
  \item[\textbf{D2}] \(\mathbb{Z}[1/2]\) \emph{is eliminated.}  The paper's offset map is
        replaced by its integer multiple \(\texttt{fnat}\,n\,a = \sum_{m<n} 3^{n-1-m}
        2^{\,a_{[1,m]}}\), so the iteration identity (1.7)\(\times 2^{|a|}\) lives entirely
        in \(\mathbb{N}\) (node~\ref{C2}).
  \item[\textbf{D3}] \emph{Asymptotics are reified.}  \(X\ll Y\) becomes
        \(\exists C>0,\ X\le C\,Y\); \(\ll_A\) becomes \(\forall A>0,\exists C>0\).  No
        \texttt{IsBigO}/filters in load-bearing statements (uniformity in \(n,\xi\) is
        exactly what filter-\(O\) would hide).
  \item[\textbf{D4}] \emph{The small constant is a numeral.}  \(\varepsilon := 10^{-4}\)
        (candidate; validated per-site).
  \item[\textbf{D6}] \emph{The infinite renewal process is finitized} by the downward
        recursion \(Q\) (node~\ref{X4}) --- the single most important de-risking decision.
\end{itemize}
Throughout, tuples are \(0\)-indexed (\(a:\mathrm{Fin}\,n\to\mathbb{N}\); paper index \(i\)
is Lean index \(i-1\)), and \(a_{[1,m]}\) denotes the prefix sum \(\sum_{i<m} a_i\).


\chapter{The support layer}\label{chp:support}

Discrete-probability, generating-function, and Fourier scaffolding used everywhere below.

\section{PMF calculus}

\begin{definition}[Discrete-probability calculus]\label{S1}
  \lean{PMF.dTV, PMF.expect, PMF.cexpect, PMF.iid}
  \leanok
Total variation distance \(d_{\mathrm{TV}}(p,q) := \sum_a' |(p\,a).\mathrm{toReal} - (q\,a).\mathrm{toReal}|\),
real and complex expectations \(\mathbb{E}_p[f] := \sum_a' (p\,a).\mathrm{toReal}\cdot f(a)\),
and the \(n\)-fold i.i.d.\ product \(p^{\otimes n}\) on \(\mathrm{Fin}\,n\to\alpha\)
(node D1).  Together with the coupling inequality (paper (1.10)), stated in indicator-expectation
form to avoid \texttt{OuterMeasure} friction: for any event \(E\subseteq\alpha\),
\[
  \bigl|\mathbb{E}_p[\mathbf 1_E] - \mathbb{E}_q[\mathbf 1_E]\bigr| \le d_{\mathrm{TV}}(p,q).
\]
Basic algebra of \(d_{\mathrm{TV}}\) (symmetry, non-negativity) and the support/apply lemmas of
the i.i.d.\ product are part of this node.
\lapsrisk{8--15}{low}{90\%}
\end{definition}

\section{The geometric family and local bounds}

\begin{definition}[Geometric and Pascal laws; negative binomial]\label{S2}
  \lean{geomHalf, geomQuarter, pascal, pascalNe3}
  \uses{S1}
  \leanok
The lattice laws driving the Syracuse valuation: \(\mathrm{Geom}(2)\) with
\(\mathbb{P}(a)=2^{-a}\) for \(a\ge 1\) (\texttt{geomHalf}); \(\mathrm{Geom}(4)\)
(\texttt{geomQuarter}); \(\mathrm{Pascal}\), the law of \(a_1+a_2\) with
\(a_i\sim\mathrm{Geom}(2)\), i.e.\ \(\mathbb{P}(b)=(b-1)2^{-b}\) for \(b\ge 2\)
(\texttt{pascal}); and its conditioning off \(b=3\) (\texttt{pascalNe3}).  The load-bearing
computation is the exact \textbf{negative-binomial point mass} for the sum of \(n\) i.i.d.\
\(\mathrm{Geom}(2)\) variables,
\[
  \mathbb{P}\!\Bigl(\textstyle\sum_{i<n} a_i = L\Bigr)
    = \binom{L-1}{n-1}\,2^{-L}, \qquad L\ge n\ge 1
\]
(compositions of \(L\) into \(n\) positive parts), together with the moment generating
functions and the one-dimensional Chernoff upper-tail bounds derived from them.
\lapsrisk{6--12}{low}{90\%}
\end{definition}

\begin{lemma}[Local 2-D Gaussian-type bound; Lemma 2.2]\label{S3}
  \lean{Gweight, geomHalf_local_bound, geomQuarter_local_bound, pascal_local_bound, hold_local_bound, hold_tail_bound}
  \uses{S2}
  \leanok
\textbf{Risk kernel 1.}  Lemma~2.2(i)(ii): local-limit / Gaussian-type \emph{upper} bounds on
the point masses of sums of \(k\) i.i.d.\ copies of the relevant lattice distributions
(\(\mathrm{Geom}(2)\), \(\mathrm{Geom}(4)\), \(\mathrm{Pascal}\), and the holding law
\(\mathrm{Hold}\) of §7), uniform in the shift, of the shape
\(\mathbb{P}(S_k = v)\ll k^{-d/2}\) with \(d\) the lattice dimension; together with the weight
function \(G_n\) (paper (2.2)).  Proved not by the paper's contour/Morera route but by explicit
generating functions plus real-variable exponential tilting (design decision D5): for the
\(d=2\) holding sums this is the circle-method form
\(\mathbb{P}(S_k=v)=(2\pi)^{-2}M(\lambda)^k e^{-\lambda\cdot v}\int_{[-\pi,\pi]^2}
|\varphi_\lambda(t)|^k\,dt\) with \(\varphi_\lambda,M\) explicit rational functions.
Statements judge-ratified 2026-07-12 vs pp.14--15: per-instance forms in
\texttt{Prob/LocalBound.lean} (scalar) and \texttt{Sec7/Unroll.lean} (\(d=2\) Hold, with
the correct \((n+1)^{-1}\) prefactor and Euclidean norm fed to the scalar \(G\)); the
\(G_{1+n}\)-for-\(G_n\) index is constants-equivalent (in-paper precedent: Lemma 7.7).
\lapsrisk{12--25}{medium --- risk kernel 1}{78\% (2026-07-12 judge re-rate: the analytic engine is largely PROVED --- circle-method inversion, character decay, center-regime local bound, generic tilting; what remains is regime assembly into the six ratified instance statements)}
\end{lemma}

\section{Fourier analysis on \(\mathbb{Z}/3^n\mathbb{Z}\)}

\begin{definition}[Additive characters and the oscillation functional]\label{S4}
  \lean{eC, osc}
  \leanok
The additive character \(e(\theta):=\exp(2\pi i\theta)\) on \(\mathbb{Q}\), the discrete
Fourier transform and Parseval identity on \(\mathbb{Z}/3^n\mathbb{Z}\), and the
\textbf{oscillation functional} \(\mathrm{Osc}_{m,n}\) (paper (1.24)) measuring how far a
weight \(c:\mathbb{Z}/3^n\mathbb{Z}\to\mathbb{R}\) is from being constant on the fibres of the
projection to \(\mathbb{Z}/3^m\mathbb{Z}\):
\[
  \mathrm{Osc}_{m,n}(c)=\sum_{Y}\Bigl|\,c(Y)-3^{\,m-n}\!\!\sum_{Y'\mapsto \bar Y} c(Y')\,\Bigr|,
  \qquad m\le n.
\]
Includes the Remark~1.18 triangle inequality that lets \(\mathrm{Osc}\) be controlled scale by
scale.
\lapsrisk{4--8}{low}{85\%}
\end{definition}


\chapter{The Collatz--Syracuse dictionary and logarithmic density}\label{chp:dictionary}

\begin{definition}[Collatz and Syracuse maps; the reduction (1.2)]\label{C1}
  \lean{col, colMin, oddPart, syr, syrMin, colMin_eq_syrMin_oddPart}
  \leanok
The Collatz map \(\mathrm{col}(N)=3N+1\) if \(N\) is odd and \(N/2\) otherwise; its orbit
minimum \(\mathrm{Colmin}(N)=\inf_k \mathrm{col}^{[k]}(N)\).  The odd part
\(\mathrm{oddPart}(N)=N/2^{v_2(N)}\), the \textbf{Syracuse map}
\(\mathrm{syr}(N)=\mathrm{oddPart}(3N+1)\) (odd \(\to\) odd), and its orbit minimum
\(\mathrm{Syrmin}\).  The elementary reduction (paper (1.2)) that lets the entire problem be
studied on the odd numbers:
\[
  0<N \;\Longrightarrow\; \mathrm{Colmin}(N) = \mathrm{Syrmin}(\mathrm{oddPart}\,N).
\]
Also the cheap odd-preservation and positivity facts (\(\mathrm{syr}\) maps odds to odds, etc.).
\lapsrisk{3--6}{low}{95\%}
\end{definition}

\begin{lemma}[Valuation vector, integer offset, and the iteration identity (1.7)]\label{C2}
  \lean{valVec, fnat, syr_iterate_key}
  \uses{C1}
  \leanok
The \(n\)-step \textbf{valuation vector} (paper (1.8))
\(\vec a^{(n)}(N)_i = v_2\bigl(3\,\mathrm{syr}^{[i]}(N)+1\bigr)\), the integerified offset
(design D2) \(\texttt{fnat}\,n\,a = \sum_{m<n} 3^{\,n-1-m}\,2^{\,a_{[1,m]}}\), and the
load-bearing \textbf{iteration identity}, the paper's (1.7) multiplied through by \(2^{|a|}\)
so it holds in \(\mathbb{N}\): for odd \(N\), with \(a=\vec a^{(n)}(N)\),
\[
  2^{\,a_{[1,n]}}\cdot \mathrm{syr}^{[n]}(N) = 3^n N + \texttt{fnat}\,n\,a .
\]
(By induction on \(n\); the step is \(2^{a_{n+1}}\mathrm{syr}(\mathrm{syr}^{[n]}N)
= 3\,\mathrm{syr}^{[n]}N + 1\).)  Together with Lemma~2.1: the valuation vector is the
\emph{unique} tuple with all entries \(\ge 1\) making the affine quotient odd.
% TODO-PRECISE: exact uniqueness shape (RATIFY-2 in the spec) still open.
\lapsrisk{5--10}{low}{90\%}
\end{lemma}

\begin{definition}[Logarithmic density and the reduction Thm 1.6 \(\Rightarrow\) Thm 1.3]\label{C3}
  \lean{logProb, posInterval, HasLogDensity, AlmostAllPos, AlmostAllOdd}
  \uses{C1}
  \leanok
The logarithmic-density apparatus: \texttt{logSum}, \texttt{logProb}, the sampling window
\texttt{posInterval}, the predicate \texttt{HasLogDensity}, and the ``almost all''
quantifiers \texttt{AlmostAllPos} and its odd-window form \texttt{AlmostAllOdd}
(logarithmic density restricted to \(2\mathbb{N}+1\)).  With the harmonic-sum integral tests
(paper (5.25)(5.26)) and the elementary splitting reducing the even-index statement
Thm~1.6 to the main theorem Thm~1.3 (node~\ref{C6}) by passing to odd parts.
\lapsrisk{6--12}{low}{85\%}
\end{definition}


\chapter{The Syracuse random variable and its valuation}\label{chp:syracuse}

\begin{definition}[The Syracuse random variable \(\mathrm{Syrac}(\mathbb{Z}/3^n\mathbb{Z})\)]\label{C4}
  \lean{syracZ, syracZ_map_cast}
  \uses{C2, S1}
  \leanok
The law \texttt{syracZ}\,\(n\) on \(\mathbb{Z}/3^n\mathbb{Z}\) (paper (1.21)/(1.26), in the
\emph{reversed} variable order deliberately, per footnote~6): the pushforward of
\(\mathrm{Geom}(2)^{\otimes n}\) under
\[
  a \;\longmapsto\; \sum_{j<n} 3^{\,j}\,(2^{-1})^{\,a_{[1,j+1]}} \in \mathbb{Z}/3^n\mathbb{Z},
\]
where \(2^{-1}\) is the inverse of the unit \(2\) mod \(3^n\).  With the projection
compatibility (paper (1.22)): \texttt{syracZ}\,\(n\) pushes forward to \texttt{syracZ}\,\(k\)
under \(\mathbb{Z}/3^n\mathbb{Z}\to\mathbb{Z}/3^k\mathbb{Z}\) for \(k\le n\); and the recursion
Lemma~1.12 (Euler's theorem plus a geometric-series normalization
\((1-2^{-2\cdot 3^n})^{-1}\)).
\lapsrisk{5--10}{low}{85\%}
\end{definition}

\begin{proposition}[Valuation \(\approx\) geometric; Prop 1.9 and the §4 tail bound]\label{C5}
  \lean{valuation_dist}
  \uses{C2, S2}
  \leanok
\textbf{Proposition 1.9.}  There are \(c_1,C>0\) such that for all \(n,n'\) and every law
\(X\) supported on odd numbers that is \(\ll 2^{-n'}\)-close in total variation to uniform
modulo \(2^{n'}\), provided \((2+c_0)n\le n'\),
\[
  d_{\mathrm{TV}}\!\bigl(X_*\vec a^{(n)},\ \mathrm{Geom}(2)^{\otimes n}\bigr)
    \le C\,2^{-c_1 n}.
\]
I.e.\ the Syracuse valuation vector of an equidistributed-mod-\(2^{n'}\) input is
exponentially close to \(n\) independent geometric variables.  Includes the §4 tail bound
Lemma~4.1 controlling \(\mathbb{P}(a_{[1,n]}\ge n')\).
\lapsrisk{10--18}{medium}{80\%}
\end{proposition}


\chapter{First passage, stabilization, and the main theorem}\label{chp:mainthm}

\section{The main theorem}

\begin{theorem}[Main theorem: Thm 1.3, Thm 1.6, quantitative Thm 3.1]\label{C6}
  \lean{tao_collatz, tao_collatz_quantitative}
  \uses{C3, C7, C8, C9}
  \leanok
\textbf{Theorem 1.3} (the summit): for every \(f:\mathbb{N}\to\mathbb{R}\) with
\(f(N)\to\infty\), \(\mathrm{Colmin}(N)<f(N)\) for almost all \(N\) in logarithmic density.
It follows from the intermediate \textbf{Theorem 1.6} (almost all orbits attain almost bounded
values) which in turn follows from the \textbf{quantitative Theorem 3.1}: there exist
\(c,C>0\) so that for all \(N_0,x\ge 2\),
\[
  \mathrm{logProb}_{[1,x]}\bigl(\{N : \mathrm{Colmin}(N)\le N_0\}\bigr)
    \;\ge\; 1 - C/(\log N_0)^{c}.
\]
Thm~3.1 is obtained from the stabilization estimate Prop~1.11 (node~\ref{C9}) by iterating
over dyadic scales (§3).
\lapsrisk{8--15}{low}{85\%}
\end{theorem}

\section{First passage (§5)}

\begin{lemma}[First-passage non-escape; (1.19)]\label{C7}
  \lean{passes, passTime, passLoc}
  \uses{C5, S2}
  \leanok
The first-passage data: \(\mathrm{passes}(x,N)\) (\(=T_x(N)<\infty\), the orbit reaches
\(\le x\)), the passage time \(\mathrm{passTime}\), and the passage location
\(\mathrm{Pass}_x = \mathrm{passLoc}\) (with the convention \(\mathrm{Syr}^\infty := 1\)).
Estimate (paper (1.19)): the probability that a log-uniformly chosen odd \(N_y\) in a window
around \(y\) \emph{never} descends to \(\le x\) is small,
\[
  \mathbb{P}\bigl(T_x(N_y)=\infty\bigr) \ll x^{-c}.
\]
\lapsrisk{5--10}{low}{85\%}
\end{lemma}

\begin{proposition}[Approximate first-passage formula; Prop 5.2]\label{C8}
  \uses{C2, C5, C7}
  \notready
\textbf{Proposition 5.2}, the approximate formula (paper (5.8)) for the law of the
first-passage location, built from the bookkeeping events
\(\mathcal{A}^{(n')}\) (5.11), \(E'\) (5.10), and \(I_y\) (5.9), through the
\(B_{n,y}\) equivalence chain relating the passage location to the affine images of the input
window.  This is where the Prop~1.9 geometric approximation is fed into the concrete window
statistics.
% TODO-PRECISE: Prop 5.2 has no canonical Lean name in SKELETON-SPEC yet; bound when named.
\lapsrisk{15--30}{medium}{75\%}
\end{proposition}

\begin{proposition}[Stabilization across scales; Prop 1.11]\label{C9}
  \lean{stabilization}
  \uses{C8, C10}
  \leanok
\textbf{Proposition 1.11}, the spine's key input.  With \(\alpha=1.001\) (paper (1.18)):
there are \(c,C,x_0>0\) so that for \(x\ge x_0\), the first-passage location laws seen from two
successive scales \(x^{\alpha}\) and \(x^{\alpha^2}\) agree up to
\[
  d_{\mathrm{TV}}\Bigl(
      \bigl(\mathrm{logUnifOdd}_{[x^{\alpha},\,x^{\alpha^2}]}\bigr)_*\mathrm{Pass}_{\lfloor x\rfloor},\
      \bigl(\mathrm{logUnifOdd}_{[x^{\alpha^2},\,x^{\alpha^3}]}\bigr)_*\mathrm{Pass}_{\lfloor x\rfloor}
    \Bigr) \le C(\log x)^{-c},
\]
together with the non-escape control on each window.  Proved via Lemma~5.3
(\(c_n(X)\ll 1\)) and (5.18)--(5.21), applying the fine-scale mixing Prop~1.14
(node~\ref{C10}) at a base scale \(m_0\).
% TODO-PRECISE: window endpoints (nested \(\alpha\)-powers, RATIFY-3) to be ratified.
\lapsrisk{10--20}{medium}{75\%}
\end{proposition}


\chapter{Fine-scale mixing (§6)}\label{chp:mixing}

\begin{proposition}[Fine-scale mixing; Prop 1.14 \(\Leftarrow\) Prop 1.17]\label{C10}
  \lean{fine_scale_mixing}
  \uses{C4, S2, S4, X11}
  \leanok
\textbf{Proposition 1.14}: the Syracuse law is asymptotically equidistributed at fine scales,
\[
  \forall A>0,\ \exists C>0,\ \forall\, 1\le m\le n:\quad
  \mathrm{Osc}_{m,n}\bigl(Y\mapsto (\texttt{syracZ}\,n)(Y)\bigr) \le C\,m^{-A}.
\]
Deduced from the \textbf{character-sum decay} Prop~1.17 (node~\ref{X11}, the head of the §7
crux ``X-chain'') by Plancherel on \(\mathbb{Z}/3^n\mathbb{Z}\).  The §6 machinery around it:
Lemma~6.2 (\(F_n\) injective), Cor~6.3 (\(3\)-adic separation of the offsets), the event \(E\)
(6.2), the stopping time \(k\), and the Plancherel step.
\lapsrisk{15--30}{medium}{75\%}
\end{proposition}


\chapter{The crux: character decay via renewal versus triangles (§7)}\label{chp:crux}

The entire descent bottoms out at Prop~1.17.  §7 proves it by studying the phase functions
\(\theta(j,l)\) attached to the paired valuations, splitting the index set into ``white''
points (where the character sum contracts) and ``black'' points (where it does not), showing
the black set is a union of separated triangles, and running a renewal process that is forced
to hit many white points --- unless it is trapped in a large triangle, which is rare.

\section{Setup and the phase functions}

\begin{definition}[§7 setup: character, pairing, factorization]\label{X1}
  \uses{C4, S2}
  \notready
The character \(\chi\) (paper (7.1)) and the reduction of Prop~1.17 to a bound on
\(\mathbb{E}\,\chi\) through the reversed form (1.26); the \textbf{pairing}
\(b_j = a_{2j-1}+a_{2j}\sim\mathrm{Pascal}\) that groups the geometric valuations two at a
time; and the conditional \textbf{factorization} (7.4)/(7.5) of the Fourier coefficient into a
product of per-pair factors \(f\) (with the low-order piece \(g\) split off).
% TODO-PRECISE: no canonical Lean name for the §7.1 setup in SKELETON-SPEC yet.
\lapsrisk{6--12}{medium}{80\%}
\end{definition}

\begin{lemma}[Phase identities and white-point cancellation; Lemma 7.2]\label{X2}
  \lean{sfrac, θq, θq_succ_j, θq_pred_l, black}
  \uses{X1}
  \leanok
The phase \(\theta(j,l)\) (paper (7.7)/(7.8)), a signed fractional part of
\(\xi\cdot 3^{2j-2}\cdot 2^{\,1-l}/3^n\), and its two scaling identities (7.13)/(7.14): there
are integers \(k\) with
\[
  \theta(j{+}1,l) = 9\,\theta(j,l) + k, \qquad \theta(j,l{-}1) = 2\,\theta(j,l) + k .
\]
The black/white dichotomy (7.9) with \(\varepsilon = 10^{-4}\): a point is \emph{black} when
\(|\theta(j,l)|\le\varepsilon\), \emph{white} otherwise.  \textbf{Lemma 7.2} (white-point
cancellation): at a white point the per-pair factor contracts,
\(|f| \le |\cos(\pi\theta)| \le \exp(-\varepsilon^3)\).
\lapsrisk{4--8}{low}{90\%}
\end{lemma}

\section{The triangle structure of the black set}

\begin{lemma}[Black set = separated triangles; Lemma 7.4]\label{X3}
  \lean{triangle, black_structure}
  \uses{X2}
  \leanok
\textbf{Risk kernel 2.}  With a triangle (paper (7.11)) the set
\(\{(j,l) : j_0\le j,\ l\le l_0,\ (j-j_0)\log 9 + (l_0-l)\log 2 \le s\}\), Lemma~7.4 asserts
the black set inside the strip is a disjoint union of triangles that are (i) confined to the
strip \(j+1\le \lfloor n/2\rfloor - \tfrac{1}{10}\log(1/\varepsilon)\), and (ii) pairwise
Euclidean-separated by at least \(\tfrac{1}{10}\log(1/\varepsilon)\):
\[
  \{(j,l) : j{+}1\le \lfloor n/2\rfloor,\ \text{black}(j,l)\}
    \;=\; \bigsqcup_{t\in T} \mathrm{triangle}(t).
\]
Proved through the phase identities (7.12)--(7.15), the weakly-black claims (i)--(iii), the
\(l^*/j^*\) construction, and Claim~(\(*\)) Cases~1--3.
% TODO-PRECISE: exact conjunction spelling (RATIFY-5) is implementer's choice.
% Risk kernel 2 CLOSED 2026-07-10 (box laps 6--10, est. was 15--30 laps / 75%).
\end{lemma}
\begin{proof}
  \leanok
  \texttt{black\_structure} in \texttt{Sec7/Triangles.lean}, via the exact fibre identity
  \(\theta(j,l) = 9^{j-j^*}2^{l^*-l}\theta^*\) (sharper than the paper's (7.18) inequality;
  at \(\varepsilon = 10^{-4}\) the separation conjunct reduces to lattice disjointness).
  Judge-verified \texttt{\#print axioms}: \texttt{[propext, Classical.choice, Quot.sound]}
  (2026-07-10, host \texttt{lake env lean}).
\end{proof}

\section{The holding process and its finitized renewal recursion (D6)}

\begin{proposition}[Holding law, the \(Q\) recursion, and Prop 7.3]\label{X4}
  \lean{hold, Q_rec, Q_boundary, renewal_white_encounters}
  \uses{S2, X2}
  \leanok
The \textbf{holding law} \(\mathrm{Hold}\) on \(\mathbb{N}\times\mathbb{Z}\) (§7.3), whose first
coordinate is always \(\ge 1\), governing how the renewal walk steps between pairs.  Design
decision D6 replaces the paper's infinite renewal product (7.34) by the \emph{defined}
downward recursion \(Q:\mathbb{N}\times\mathbb{Z}\to\mathbb{R}\) (paper (7.35)),
\[
  Q(j,l) = 1 \ \ (j>\lfloor n/2\rfloor); \qquad
  Q(j,l) = e^{-\varepsilon^3\mathbf 1_W(j,l)}\cdot\!\!\sum_{d}\!(\mathrm{Hold}\,d)\,Q\bigl((j,l)+d\bigr)
  \ \ (j\le\lfloor n/2\rfloor),
\]
well-founded on \(\lfloor n/2\rfloor - j\) (Hold's first coordinate is \(\ge 1\)), with
\(0\le Q\le 1\).  The bridge (7.28)/(7.34)--(7.36) makes
\(\mathbb{E}\,Q(\mathrm{Hold})\ll_A n^{-A}\) equivalent to \textbf{Proposition 7.3}
(the renewal white-encounter count, in the finite-vector form over \(\mathrm{Pascal}^{\otimes\lfloor n/2\rfloor}\)).
This node exists to validate the D6 finitization.
\lapsrisk{8--15}{medium}{80\%}
\end{proposition}

\begin{lemma}[Holding-law basics; Lemma 7.6]\label{X5}
  \uses{X4}
  \notready
\textbf{Lemma 7.6}: the explicit distribution of \(\mathrm{Hold}\), its exponential tail,
aperiodicity, and mean vector \((4,16)\).  These are the renewal-theory inputs (drift and
spread) for the first-passage analysis of the walk.
% TODO-PRECISE: no canonical Lean name in SKELETON-SPEC yet.
\lapsrisk{4--8}{low}{85\%}
\end{lemma}

\begin{lemma}[First-passage location of the holding walk; Lemma 7.7]\label{X6}
  \lean{fpDist, fpDist_location_bound}
  \uses{S3, X5}
  \leanok
\textbf{Lemma 7.7} (paper p.43): the first-passage endpoint mass satisfies
\(\mathbb{P}(\mathbf{v}_{[1,\mathbf{k}]}=(j,l)) \ll e^{-c(l-s)}(1+s)^{-1/2}
G_{1+s}(c(j-s/4))\).  D6 form: \texttt{fpDist} (budget-recursive endpoint PMF, support
facts proved) and \texttt{fpDist\_location\_bound}, judge-ratified vs the p.43 display
2026-07-10 (unconditional in \(l\) since the support fact kills \(l\le s\)).
Uses the local 2-D bound (node~\ref{S3}) applied to the holding sums.
\lapsrisk{10--20}{high}{70\%}
\end{lemma}

\section{The three cases of Prop 7.8}

\begin{proposition}[Monotone functional \(Q_m\); Prop 7.8 skeleton; Case 1]\label{X7}
  \lean{Qm, prop_7_8, Q_white_case1, Q_polynomial_decay}
  \uses{X4}
  \leanok
The truncated functional \(Q_m\) (paper (7.38)) and the skeleton of \textbf{Proposition 7.8}
(the quantitative white-encounter bound), together with \textbf{Case 1} (paper (7.42)/(7.43)):
when the walk starts at a white point, the \(e^{-\varepsilon^3}\) contraction from Lemma~7.2
(node~\ref{X2}) already delivers the required decay.  Judge-ratified 2026-07-10 vs pp.45--46:
(7.38)\(\to\)\texttt{Qm}, (7.40)\(\to\)\texttt{prop\_7\_8}, (7.43)\(\to\)\texttt{Q\_white\_case1}
(verbatim incl.\ \(e^{-\varepsilon^3/2}\)), (7.37)\(\to\)\texttt{Q\_polynomial\_decay}.  Case 1
is \emph{proved}; the skeleton delegates the black edge to \texttt{Q\_black\_edge}
(nodes~\ref{X8}/\ref{X10}).
\lapsrisk{4--8}{low}{85\%}
\end{proposition}

\begin{lemma}[Case 2: shallow inside a triangle]\label{X8}
  \uses{X3, X6, X7}
  \notready
\textbf{Risk kernel 3a.}  \textbf{Case 2} (paper (7.44)--(7.51)): the walk starts shallow
inside a triangle; the analysis produces the \(\gg 1\) white-exit bound (7.50)/(7.51) showing
the walk leaves the triangle through many white points with high probability.  Combines the
triangle geometry (node~\ref{X3}), the first-passage location law (node~\ref{X6}), and the
Case-1 contraction (node~\ref{X7}).
% TODO-PRECISE: no canonical Lean name in SKELETON-SPEC yet.
\lapsrisk{12--25}{high --- risk kernel 3a}{70\% (2026-07-10: white-exit probability $\approx 0.99$ Monte Carlo; $\varepsilon$-sites verified)}
\end{lemma}

\begin{lemma}[Many triangles \(\Rightarrow\) many white points; Lemma 7.9]\label{X9}
  \uses{X4, X8}
  \notready
\textbf{Lemma 7.9}: an induction on \(R\) over the \(Q\)-recursion (node~\ref{X4}) showing that
a walk crossing many triangles must encounter many white points, so the accumulated
\(e^{-\varepsilon^3}\) factors give the exponential decay.
% TODO-PRECISE: no canonical Lean name in SKELETON-SPEC yet.
\lapsrisk{8--15}{high}{70\%}
\end{lemma}

\begin{lemma}[Large triangles rarely encountered; Lemma 7.10]\label{X10}
  \uses{X3, X6}
  \notready
\textbf{Risk kernel 3b.}  \textbf{Lemma 7.10} (paper (7.60)--(7.65)): after a lengthy crossing,
large triangles are rarely encountered, by a separated-\(\Sigma\) counting argument over the
triangle family (node~\ref{X3}) using the first-passage location law (node~\ref{X6}).  This is
the complementary tail to Case~2.
% TODO-PRECISE: no canonical Lean name in SKELETON-SPEC yet.
\lapsrisk{15--30}{high --- risk kernel 3b}{65\% (2026-07-10: row-disjointness + $\Sigma$ $j$-separation verified on real triangle inventory)}
\end{lemma}

\section{Assembly: Prop 1.17}

\begin{proposition}[Case 3 assembly; Prop 7.8 \(\to\) 7.3 \(\to\) 7.1 \(\to\) Prop 1.17]\label{X11}
  \lean{charFn_decay, key_fourier_decay}
  \uses{X9, X10}
  \leanok
\textbf{Case 3} assembly (paper (7.48)/(7.49), (7.52)--(7.67)): the events \(E_*,F_*\), the
choice \(R=\lfloor A^2/\varepsilon^4\rfloor\), and the deterministic claim (7.67), combining
Lemma~7.9 (node~\ref{X9}) and Lemma~7.10 (node~\ref{X10}).  This closes Prop~7.8, hence
Prop~7.3 (node~\ref{X4}) and Prop~7.1, hence the \textbf{crux} \textbf{Proposition 1.17}
(character-sum decay): for every \(A>0\) there is \(C>0\) with, for all \(n\ge 1\) and all
\(\xi\) with \(3\nmid\xi\),
\[
  \Bigl\|\,\mathbb{E}_{Y\sim\texttt{syracZ}\,n}\, e\bigl(-\xi\, Y/3^{n}\bigr)\Bigr\|
    \;\le\; C\,n^{-A}.
\]
This is the single Fourier-analytic input the whole spine reduces to (feeds node~\ref{C10}).
\lapsrisk{10--20}{high}{70\%}
\end{proposition}
